% Generated draft: requires mathematical and source review.
\documentclass[11pt]{article}
\usepackage[margin=25mm]{geometry}
\usepackage{fontspec}
\setmainfont{Latin Modern Roman}
\usepackage{amsmath,amssymb,mathtools,unicode-math}
\setmathfont{Latin Modern Math}
\usepackage{enumitem,needspace,xurl,hyperref}
\setlength{\parindent}{0pt}
\setlength{\parskip}{5pt}
\setlength{\emergencystretch}{3em}
\newcommand{\N}{\mathbb{N}}
\newcommand{\Z}{\mathbb{Z}}
\newcommand{\Q}{\mathbb{Q}}
\newcommand{\R}{\mathbb{R}}
\newcommand{\C}{\mathbb{C}}
\newcommand{\F}{\mathbb{F}}
\newcommand{\A}{\mathbb{A}}
\newcommand{\PP}{\mathbb{P}}
\newcommand{\E}{\mathbb{E}}
\newcommand{\eps}{\varepsilon}
\newcommand{\dd}{\,\mathrm{d}}
\begin{document}
\section*{UnsolvedMath ID: 800048}
\textbf{Source problem code:} catalogue entry
\subsection{Definitions and mathematical statement}
Let a semigroup be a set equipped with an associative binary operation. For a semigroup \(S\), write
\[
E(S)=\{e\in S:e^2=e\}.
\]
For a set \(X\), let \(T_X\) be the semigroup of all functions \(X\to X\) under composition. In the source's terminology, a projection is an idempotent transformation, so the projections in \(T_X\) are precisely
\[
E(T_X)=\{p:X\to X:p\circ p=p\}.
\]
No linear structure or orthogonality condition is specified.

\textbf{Requested characterizations (E.\,S. Lyapin, Problem 50).}
\begin{enumerate}
\item Characterize the pargoids of projections in semigroups of all transformations.
\item Characterize the pargoids of matrix units.
\item Characterize the pargoids of idempotents of the basic classes of semigroups.
\end{enumerate}

\textbf{Unresolved formulation qualifications.} The supplied record does not define a pargoid or explain how its underlying set and partial multiplication are obtained from the indicated semigroup data. Consequently, the domains of the partial operations and the precise objects to be characterized cannot be specified from this record alone. In particular, no rule declaring a product defined exactly when it remains in a selected subset is assumed. The record also does not specify restrictions on the underlying transformation sets, the index sets or zero convention for matrix units, or which semigroup classes are meant by ``the basic classes.'' These choices require clarification before the three requests can be made fully formal.

The source reports that some classes of pargoids were characterized in the paper labelled [ESL5], without identifying those classes or supplying the paper's bibliographic details.
\subsection{Short English statement}
Problem 50, attributed to E.\,S. Lyapin, asks for characterizations of pargoids arising from projections, matrix units, and idempotents in basic classes of semigroups. Here projections mean idempotent transformations. The supplied record leaves the pargoid construction and several target classes unspecified, so this is a statement-only draft requiring human mathematical review.
\subsection{Sources}
\begin{itemize}
\item Supplied local registry record: permanent UnsolvedMath ID \(800048\), source problem code \texttt{catalogue entry}, title \emph{The Lyapin's Notebook --- Problem 50}; collection: AMR Open Problem Lists. The registry labels the problem \texttt{open}; this is a supplied status, not an independently verified mathematical assessment.
\item Bershadsky, Kublanovsky, and Mashevitzky, \emph{The Lyapin's Notebook: Unsolved Problems in Semigroup Theory} (2026), Problem 50. Supplied source URL: \url{https://arxiv.org/abs/2604.04763}. The background records access on 2026-10-16 and extraction by \texttt{arxiv\_latex}; these are supplied provenance details, not a claim of independent access or verification.
\item The problem text attributes the questions to E.\,S. Lyapin. The registry's \texttt{proposed\_by} field lists Bershadsky, Kublanovsky, and Mashevitzky, with proposed year 2026; these registry fields do not resolve the original proposal date.
\item [ESL5]: cited in the supplied statement as characterizing some classes of pargoids. No title, authorship, publication details, link, or precise results for this reference are supplied.
\end{itemize}
\end{document}
